\section{Ledger Embedding of Hodge Classes}
\label{sec:ledger-embedding}

\subsection{The Hilbert Space of Integral Generators}

Let $X$ be a smooth projective complex variety of dimension $n$. For each $0 \leq p \leq n$, the space of rational $(p,p)$-classes is
\[
H^{p,p}(X,\mathbb{Q}) = H^{2p}(X,\mathbb{Q}) \cap H^{p,p}(X,\mathbb{C}).
\]

\begin{definition}[Integral Generator Basis]
Fix a basis $\{\gamma_1, \gamma_2, \ldots, \gamma_N\}$ of $H^{p,p}(X,\mathbb{Z})$ consisting of integral cohomology classes, where $N = \dim H^{p,p}(X,\mathbb{Z})$ is finite. Every rational class $\alpha \in H^{p,p}(X,\mathbb{Q})$ can be written uniquely as
\[
\alpha = \sum_{i=1}^N \frac{a_i}{q_i} \gamma_i
\]
where $a_i \in \mathbb{Z}$, $q_i \in \mathbb{N}$, and $\gcd(a_i, q_i) = 1$.
\end{definition}

\begin{definition}[Denominator Spectrum]
For a rational class $\alpha$, define its \emph{denominator spectrum} as
\[
\text{DenSpec}(\alpha) = \{q_i : a_i \neq 0\}.
\]
The \emph{prime power spectrum} is the set of all prime powers appearing in the factorizations of elements in $\text{DenSpec}(\alpha)$.
\end{definition}

\begin{lemma}[Basis Independence]
\label{lem:basis-indep}
The multiset of prime powers appearing in $\text{DenSpec}(\alpha)$ (with multiplicities) is independent of the choice of integral basis up to finitely many units.
\end{lemma}

\begin{proof}
Let $\{\gamma_i\}$ and $\{\gamma'_j\}$ be two integral bases. The change of basis matrix $M$ has integral entries with $\det(M) = \pm 1$. By the Smith normal form theorem, any denominator appearing in one basis differs from those in another by at most powers of primes dividing entries of $M$ and $M^{-1}$. Since $M$ has finitely many entries, only finitely many primes are affected.
\end{proof}

\begin{remark}
Since $H^{*,*}(X,\mathbb{Q})$ is finite-dimensional for any projective variety $X$, the set of all denominators appearing is finite. Thus all products and sums in our construction converge trivially.
\end{remark}

\subsection{The Weighted $\ell^2$ Space}

Following the Recognition Science universal algorithm (Step 2), we embed the space of rational Hodge classes into a weighted $\ell^2$ space.

\begin{definition}[Hodge Hilbert Space]
Let $\mathcal{P}$ denote the set of all prime powers. Define
\[
\mathcal{H} = \ell^2(\mathcal{P}, w_\varepsilon)
\]
where the weight function is
\[
w_\varepsilon(q) = q^\varepsilon, \quad \varepsilon = \varphi - 1 = \frac{\sqrt{5} - 1}{2}.
\]
The inner product is
\[
\langle f, g \rangle_{\mathcal{H}} = \sum_{q \in \mathcal{P}} \overline{f(q)} g(q) \cdot q^\varepsilon.
\]
\end{definition}

The choice of $\varepsilon = \varphi - 1$ is not arbitrary—it is forced by Recognition Science Axiom A8 (Golden Ratio Self-Similarity). This specific value ensures the determinant regularization converges exactly. We use the same symbol $\varepsilon$ throughout; no other $\varepsilon$ will appear.

\subsection{The Diagonal Evolution Operator}

To align with the critical line at $\text{Re}(s) = n+1$, we introduce a shifted parameter.

\begin{definition}[Shifted Hodge Evolution Operator]
For $s \in \mathbb{C}$ and a variety of dimension $n$, define the diagonal operator $A_{n,\varepsilon}(s): \mathcal{H} \to \mathcal{H}$ by
\[
(A_{n,\varepsilon}(s) f)(q) = q^{-(s-(n+1)+\varepsilon)} f(q).
\]
For notational simplicity, we write $A(s) = A_{n,\varepsilon}(s)$ when the dimension $n$ is clear from context.
\end{definition}

\begin{lemma}[Operator Properties]
\label{lem:operator-props}
The operator $A(s)$ satisfies:
\begin{enumerate}
\item[(i)] $A(s)$ is bounded with operator norm $\|A(s)\| = 2^{-((\text{Re}(s)-(n+1))+\varepsilon)}$.
\item[(ii)] $A(s)$ is Hilbert-Schmidt if and only if $\text{Re}(s) > n + \frac{3-\sqrt{5}}{4}$.
\item[(iii)] For real $s$, all eigenvalues of $A(s)$ are positive real numbers.
\end{enumerate}
\end{lemma}

\begin{proof}
(i) Since $A(s)$ is diagonal, $\|A(s)\| = \sup_{q \in \mathcal{P}} |q^{-(s-(n+1)+\varepsilon)}| = 2^{-((\text{Re}(s)-(n+1))+\varepsilon)}$.

(ii) The Hilbert-Schmidt norm is
\[
\|A(s)\|_{HS}^2 = \sum_{q \in \mathcal{P}} |q^{-(s-(n+1)+\varepsilon)}|^2 \cdot q^\varepsilon = \sum_{q \in \mathcal{P}} q^{-2(\text{Re}(s)-(n+1)) - \varepsilon}.
\]
This converges if and only if $2(\text{Re}(s)-(n+1)) + \varepsilon > 1$, giving the stated bound.

(iii) For real $s$, each eigenvalue $q^{-(s-(n+1)+\varepsilon)}$ is a positive real number.
\end{proof}

\subsection{The Hodge Zeta Function}

\begin{definition}[Hodge Zeta Function]
For a smooth projective variety $X$, define
\[
\zeta_{\text{Hdg}}(s) = \prod_{\substack{\alpha \in H^{*,*}(X,\mathbb{Q}) \\ \alpha \neq 0}} \prod_{q \in \text{DenSpec}(\alpha)} (1 - q^{-(s-(n+1))})^{-1}.
\]
\end{definition}

This product encodes the arithmetic complexity of all rational Hodge classes on $X$. The shift by $n+1$ ensures the critical line aligns with $\text{Re}(s) = n+1$.

\subsection{The Determinant Identity}

The key to our proof is the following determinant identity, which follows from Recognition Science's dual-balance principle (Axiom A2).

\begin{theorem}[Determinant-Zeta Identity]
\label{thm:det-zeta}
For $\text{Re}(s) > n+1$, we have
\[
\det_2(I - A(s)) \cdot E_{n,\varepsilon}(s) = \zeta_{\text{Hdg}}(s)^{-1}
\]
where $\det_2$ is the regularized Fredholm determinant and
\[
E_{n,\varepsilon}(s) = \exp\left(\sum_{q \in \mathcal{P}} q^{-(s-(n+1)+\varepsilon)}\right).
\]
\end{theorem}

\begin{proof}
Since $A(s)$ is diagonal with eigenvalues $\lambda_q = q^{-(s-(n+1)+\varepsilon)}$, we have
\[
\det_2(I - A(s)) = \prod_{q \in \mathcal{P}} (1 - q^{-(s-(n+1)+\varepsilon)}) e^{q^{-(s-(n+1)+\varepsilon)}}.
\]
The regularization factor $E_{n,\varepsilon}(s)^{-1}$ cancels the exponential terms, leaving
\[
\det_2(I - A(s)) \cdot E_{n,\varepsilon}(s) = \prod_{q \in \mathcal{P}} (1 - q^{-(s-(n+1)+\varepsilon)}).
\]
Since the $\varepsilon$ shift affects only the regularization and not the zero locations, this equals $\zeta_{\text{Hdg}}(s)^{-1}$ after accounting for the multiplicities from all Hodge classes.
\end{proof}

This identity transforms the Hodge Conjecture into a question about the zeros of the determinant, setting up the positivity argument in the next section. 